\documentclass[10pt,a4paper]{amsart}
\usepackage{amsmath,amssymb,amsthm,graphicx}
% Retained minimal preamble for the Stacks Project chapter flat.tex (More on Flatness).
% Verbatim excerpt of the chapter preamble: Xy-pic setup (native preamble lines 14-18) and the native
% theorem declarations with the mathematical macros (native preamble lines 154-217). Copied mechanically;
% nothing added, removed or redefined. The wrapper class and the amsmath/amssymb/amsthm/graphicx base
% packages are supplied by the assembler, so the class/externaldocument part of the native preamble is omitted.
% For commutative diagrams we use Xy-pic
\usepackage[all]{xy}

% We use 2cell for 2-commutative diagrams.
\xyoption{2cell}
% Theorem environments.
%
\theoremstyle{plain}
\newtheorem{theorem}[subsection]{Theorem}
\newtheorem{proposition}[subsection]{Proposition}
\newtheorem{lemma}[subsection]{Lemma}

\theoremstyle{definition}
\newtheorem{definition}[subsection]{Definition}
\newtheorem{example}[subsection]{Example}
\newtheorem{exercise}[subsection]{Exercise}
\newtheorem{situation}[subsection]{Situation}

\theoremstyle{remark}
\newtheorem{remark}[subsection]{Remark}
\newtheorem{remarks}[subsection]{Remarks}

\numberwithin{equation}{subsection}

% Macros
%
\def\lim{\mathop{\mathrm{lim}}\nolimits}
\def\colim{\mathop{\mathrm{colim}}\nolimits}
\def\Spec{\mathop{\mathrm{Spec}}}
\def\Hom{\mathop{\mathrm{Hom}}\nolimits}
\def\Ext{\mathop{\mathrm{Ext}}\nolimits}
\def\SheafHom{\mathop{\mathcal{H}\!\mathit{om}}\nolimits}
\def\SheafExt{\mathop{\mathcal{E}\!\mathit{xt}}\nolimits}
\def\Sch{\mathit{Sch}}
\def\Mor{\mathop{\mathrm{Mor}}\nolimits}
\def\Ob{\mathop{\mathrm{Ob}}\nolimits}
\def\Sh{\mathop{\mathit{Sh}}\nolimits}
\def\NL{\mathop{N\!L}\nolimits}
\def\CH{\mathop{\mathrm{CH}}\nolimits}
\def\proetale{{pro\text{-}\acute{e}tale}}
\def\etale{{\acute{e}tale}}
\def\QCoh{\mathit{QCoh}}
\def\Ker{\mathop{\mathrm{Ker}}}
\def\Im{\mathop{\mathrm{Im}}}
\def\Coker{\mathop{\mathrm{Coker}}}
\def\Coim{\mathop{\mathrm{Coim}}}

% Boxtimes
%
\DeclareMathSymbol{\boxtimes}{\mathbin}{AMSa}{"02}

%
% Macros for moduli stacks/spaces
%
\def\QCohstack{\mathcal{QC}\!\mathit{oh}}
\def\Cohstack{\mathcal{C}\!\mathit{oh}}
\def\Spacesstack{\mathcal{S}\!\mathit{paces}}
\def\Quotfunctor{\mathrm{Quot}}
\def\Hilbfunctor{\mathrm{Hilb}}
\def\Curvesstack{\mathcal{C}\!\mathit{urves}}
\def\Polarizedstack{\mathcal{P}\!\mathit{olarized}}
\def\Complexesstack{\mathcal{C}\!\mathit{omplexes}}
% \Pic is the operator that assigns to X its picard group, usage \Pic(X)
% \Picardstack_{X/B} denotes the Picard stack of X over B
% \Picardfunctor_{X/B} denotes the Picard functor of X over B
\def\Pic{\mathop{\mathrm{Pic}}\nolimits}
\def\Picardstack{\mathcal{P}\!\mathit{ic}}
\def\Picardfunctor{\mathrm{Pic}}
\def\Deformationcategory{\mathcal{D}\!\mathit{ef}}

% --- appended 2026-10-12 (Chief): the project own definitions of the reference/slogan/history
% display environments, copied mechanically from the shared Stacks preamble lines 8-12
% (\usepackage{verbatim} plus the three \newenvironment lines, which the chapter preamble keeps
% verbatim from the project preamble). Nothing added, removed or redefined: these environments
% carry no mathematical content and the project itself typesets them as comments. Without these
% three lines a counted statement whose own text uses them fails with Environment slogan undefined.
\usepackage{verbatim}
\newenvironment{reference}{\comment}{\endcomment}
\newenvironment{slogan}{\comment}{\endcomment}
\newenvironment{history}{\comment}{\endcomment}
\usepackage{cleveref}
\crefname{originallemma}{original lemma}{original lemmas}
\Crefname{originallemma}{Original lemma}{Original lemmas}
\makeatletter
\newlabel{lemma-finite-type-flat-along-fibre}{{05KU}{source}{}{sourcepointer-1}{}}
\newlabel{lemma-finite-type-flat-along-fibre@cref}{{[originallemma][0][]05KU}{[1][0][]source}{}{}{}}
\makeatother
\crefname{originallemma}{original lemma}{original lemmas}
\Crefname{originallemma}{Original lemma}{Original lemmas}
\makeatletter
\newlabel{lemma-finite-type-flat-at-point}{{05KT}{source}{}{sourcepointer-2}{}}
\newlabel{lemma-finite-type-flat-at-point@cref}{{[originallemma][0][]05KT}{[1][0][]source}{}{}{}}
\makeatother
\crefname{originallemma}{original lemma}{original lemmas}
\Crefname{originallemma}{Original lemma}{Original lemmas}
\makeatletter
\newlabel{more-morphisms-lemma-elementary-etale-neighbourhoods}{{057B}{source}{}{sourcepointer-3}{}}
\newlabel{more-morphisms-lemma-elementary-etale-neighbourhoods@cref}{{[originallemma][0][]057B}{[1][0][]source}{}{}{}}
\makeatother
\title{A finite type module flat along a fibre is flat on an open after an elementary etale neighbourhood}
\author{The Stacks Project Authors}
\date{stacks-project:flat.tex:local-snapshot-20261011; GNU FDL 1.2 or later}
\begin{document}
\maketitle
\noindent\textit{Editorial note.} All delivered mathematical text is the author's own, copied from the pinned Stacks Project chapter file flat.tex; no statement or proof step was rewritten, reordered or removed.\par\medskip
\noindent\textit{Editorial note.} Preamble provenance: the chapter's own retained minimal preamble (sources/stacks-flat/retained-preamble.tex, lines 1-85) is a verbatim excerpt of the native Xy-pic setup and of the native theorem declarations and mathematical macros, followed by the project's own reference/slogan/history display-environment definitions copied mechanically from the shared preamble; the wrapper class and the amsmath/amssymb/amsthm/graphicx base packages are supplied by the assembler, and extra\_packages is left empty because everything needed is already in the retained preamble. No macro was invented or redefined.\par\medskip
\noindent\textit{Editorial note.} Same-chapter pointers are located at their \textbackslash{}label lines inside the pinned chapter file; cross-chapter pointers are located in the project's own tag list staged in the snapshot (sources/stacks-chow/tags.txt), at the verbatim lines carrying their canonical tags.\par\medskip
\noindent\textit{Allowed standard prerequisites (editorial).}\begin{itemize}
\item The pointwise version of the flat along the fibre statement. Same chapter flat.tex, label lemma-finite-type-flat-at-point, invoked at proof line 2819. Produces, for each point of the fibre, an elementary etale neighbourhood and an open neighbourhood of the point over the base change on which the pullback of the module is flat; assembling finitely many of them is what the counted proof supplies.
\item Existence of an elementary etale neighbourhood dominating finitely many given ones. More on Morphisms, label more-morphisms-lemma-elementary-etale-neighbourhoods, invoked at proof line 2832. Supplies one elementary etale neighbourhood of the base dominating the finitely many pointwise neighbourhoods, which is what makes the union of the pieces an open neighbourhood of the fibre over a single base neighbourhood.
\item The finitely presented version of the same statement. Same chapter flat.tex, label lemma-finite-type-flat-along-fibre, invoked at proof line 2816. The author notes in the first sentence of the proof that the only difference from that lemma is the absence of the finite presentation hypothesis on the morphism.
\end{itemize}\medskip
\ifdefined\hypertarget\hypertarget{sourcepointer-1}{}\fi
\noindent\textit{Source reference.} Original lemma 05KU: sources/stacks-flat/pinned-original-source.tex, source lines 2745--2745. Navigation pointer for \textbackslash{}ref\{lemma-finite-type-flat-along-fibre\} used inside the counted span (line 2816); the argument that consumes it is retained in the span.\par\medskip
\ifdefined\hypertarget\hypertarget{sourcepointer-2}{}\fi
\noindent\textit{Source reference.} Original lemma 05KT: sources/stacks-flat/pinned-original-source.tex, source lines 2701--2701. Navigation pointer for \textbackslash{}ref\{lemma-finite-type-flat-at-point\} used inside the counted span (line 2819); the argument that consumes it is retained in the span.\par\medskip
\ifdefined\hypertarget\hypertarget{sourcepointer-3}{}\fi
\noindent\textit{Source reference.} Original lemma 057B: sources/stacks-chow/tags.txt, source lines 6739--6739. Navigation pointer for \textbackslash{}ref\{more-morphisms-lemma-elementary-etale-neighbourhoods\} used inside the counted span (line 2832); the argument that consumes it is retained in the span.\par\medskip
\begin{lemma}
\label{lemma-finite-type-flat-along-fibre-variant}
Let $f : X \to S$ be a morphism of schemes.
Let $\mathcal{F}$ be a quasi-coherent sheaf on $X$.
Let $s \in S$.
Assume that
\begin{enumerate}
\item $f$ is of finite type,
\item $\mathcal{F}$ is of finite type, and
\item $\mathcal{F}$ is flat over $S$ at every point of the fibre $X_s$.
\end{enumerate}
Then there exists an elementary \'etale neighbourhood $(S', s') \to (S, s)$
and an open subscheme
$$
V \subset X \times_S \Spec(\mathcal{O}_{S', s'})
$$
which contains the fibre $X_s = X \times_S s'$ such that the pullback
of $\mathcal{F}$ to $V$ is flat over $\mathcal{O}_{S', s'}$.
\end{lemma}

\begin{proof}
(The only difference between this and
Lemma \ref{lemma-finite-type-flat-along-fibre}
is that we do not assume $f$ is of finite presentation.)
For every point $x \in X_s$ we can use
Lemma \ref{lemma-finite-type-flat-at-point}
to find an elementary \'etale neighbourhood $(S_x, s_x) \to (S, s)$
and an open $V_x \subset X \times_S \Spec(\mathcal{O}_{S_x, s_x})$
such that $x \in X_s = X \times_S s_x$ is contained in $V_x$ and such that
the pullback of $\mathcal{F}$ to $V_x$ is flat over
$\mathcal{O}_{S_x, s_x}$. In particular we may view the fibre
$(V_x)_{s_x}$ as an open neighbourhood of $x$ in $X_s$.
Because $X_s$ is quasi-compact we can find a finite number of points
$x_1, \ldots, x_n \in X_s$ such that $X_s$ is the union of
the $(V_{x_i})_{s_{x_i}}$. Choose an elementary \'etale neighbourhood
$(S' , s') \to (S, s)$ which dominates each of the neighbourhoods
$(S_{x_i}, s_{x_i})$, see
More on Morphisms,
Lemma \ref{more-morphisms-lemma-elementary-etale-neighbourhoods}.
Set $V = \bigcup V_i$ where $V_i$ is the inverse images of the open
$V_{x_i}$ via the morphism
$$
X \times_S \Spec(\mathcal{O}_{S', s'})
\longrightarrow
X \times_S \Spec(\mathcal{O}_{S_{x_i}, s_{x_i}})
$$
By construction $V$ contains $X_s$ and by construction the pullback
of $\mathcal{F}$ to $V$ is flat over $\mathcal{O}_{S', s'}$.
\end{proof}

\end{document}
